% Faithful Korean companion of technical-report.tex
% English source SHA-256: bec9154a3a4e0555e4e6871e8cbefe11a77d0a84906b29af791daabbd918bb84
\documentclass[11pt,a4paper]{article}

\usepackage{microtype}
\usepackage{booktabs}
\usepackage{float}
\usepackage{hyperref}
\usepackage{fontspec}
\usepackage{xeCJK}
\usepackage[a4paper,margin=24mm]{geometry}
\xeCJKsetup{CJKspace=true}
\setCJKmainfont[BoldFont={NanumMyeongjoOTF Bold}]{NanumMyeongjoOTF}
\setCJKsansfont[BoldFont={Noto Sans KR Bold}]{Noto Sans KR}
\setCJKmonofont{D2Coding}
\setlength{\parindent}{0pt}
\setlength{\parskip}{0.65em}
\linespread{1.08}
\renewcommand{\abstractname}{요약}
\renewcommand{\figurename}{그림}
\renewcommand{\tablename}{표}
\renewcommand{\refname}{참고문헌}

\hypersetup{
  hidelinks,
  pdftitle={고정된 마감 아래의 근거 기반 리뷰: ICML 형식 리뷰를 위한 신뢰성 계약},
  pdfauthor={Anonymous},
  pdfsubject={Track 2 Technical Report Korean Companion},
  pdfkeywords={scientific review, evidence trace, reliable orchestration}
}
\title{고정된 마감 아래의 근거 기반 리뷰\\
  \large ICML 형식 리뷰를 위한 신뢰성 계약\\[0.5em]
  \normalsize 이해용 한국어 해설본}
\author{익명 Track 2 프로젝트}
\date{2026년 7월 12일}

\begin{document}

\maketitle

\begin{center}
\small\emph{이 문서는 이해를 돕기 위한 한국어 동반본이다. 공식 제출 형식은
영문 ICML 보고서 \texttt{technical-report.pdf}를 기준으로 한다.}
\end{center}

\begin{abstract}
시간 제한이 있는 과학 리뷰에서는 근거의 질과 안정적인 일괄 실행이
결합된다. 본 보고서는 논문 신원을 동결하고, 플랫폼 부작용을 하나의 Root
코디네이터로 제한하며, 최대 세 개의 논문별 Review Worker를 사용하도록
설계한 익명 Track 2 시스템을 제시한다. 각 Worker는 비공개 claim map,
falsification, score anchoring, consistency pass를 수행한다. 결정론적 검증
후에는 복합 고위험 초안만 선택적으로 제한된 verifier 경로에 들어간다. 10편
규모의 결정론적 합성 모의실험 세 번에서 30개 초안 모두 mock verified
completion에 도달했고, 스키마 유효성은 100\%였으며 중복 post 시도와 중복
멱등성 키는 없었다. 장애 실행에서는 malformed output, Worker timeout,
claim timeout, post timeout을 각각 한 번 복구했다. 별도의 두 프로세스 resume
검사는 4편까지의 원장 prefix와 완료 산출물을 보존한 뒤 10/10에 도달했다.
규칙 기반 seeded-fixture 검사는 one-finding baseline의 F1 0.6667에 비해 F1
1.0을 기록했다. 이 결과는 로컬 계약과 복구 로직만 검증하며, LLM 리뷰 품질,
사람 평가와의 일치도, live-platform 처리량을 입증하지 않는다.
\end{abstract}

\section{문제와 계약}

Track 2는 ICML 형식의 리뷰 에이전트와 동결 논문에 대한 구조화 리뷰를
요구한다. live 운영 시간은 25분이므로 리뷰 품질과 일괄 실행 신뢰성은 서로
연결된다. 잘 작성된 초안도 유실되거나 중복되거나 확인 없이 post되면 쓸모가
없다. 따라서 이 시스템은 근거 신원, 리뷰 유효성, 부작용 상태를 하나의
end-to-end 계약으로 다룬다.

공식 upstream Skill은 commit \texttt{a9f4f258...153473f}에 고정되어 있다.
Track 2-only 경로에는 새로운 학습, GPU, 실험 서비스, 학습 metric이 필요하지
않다. 대신 동결된 논문과 근거 hash, 독립적으로 동결된 에이전트 정의,
구조화 결과, 제공 근거가 부족할 때의 명시적 보고가 필요하다. 로컬 wrapper는
upstream의 \textit{Contribution} 평가를 보존하면서 행사 필드인
\textit{Significance}, \textit{Originality}, \textit{Comment}를 유지하며,
서로 다른 필드를 변환하지 않는다.

\section{시스템 설계}

\subsection{플랫폼 부작용의 단일 소유자}

Root 코디네이터만 assignment discovery, claim, 원격 status 확인, review post,
공유 원장 갱신을 수행할 수 있다. live 설계는 최대 세 개의 지속형 Review
Worker를 사용한다. 각 Worker는 한 번에 하나의 동결 논문 manifest를 받고
구조화된 \texttt{ReviewDraft}만 반환한다. 이 권한 경계는 Worker 재시도가
중복 플랫폼 동작으로 이어지는 것을 막는다. 합성 모의실험 runtime은 이
제한된 queue를 결정론적 로컬 drafting으로 모델링하며, native-agent 품질
실행이 아니다.

\begin{figure}[H]
  \centering
  \fbox{\parbox{0.92\columnwidth}{
    \centering
    직렬 플랫폼 경로: discover, claim, reconcile, post, verify
    \par\smallskip
    $\downarrow$ 동결 manifest와 제한된 lease
    \par\smallskip
    최대 세 Review Worker: 논문별 근거 분석만 수행
    \par\smallskip
    $\downarrow$ \texttt{ReviewDraft}
    \par\smallskip
    결정론적 검증 $\rightarrow$ 복합 risk gate
    \par\smallskip
    fast path 또는 제한된 read-only verifier $\rightarrow$ 원장
  }}
  \caption{권한은 구조적으로 좁게 제한된다. Worker는 원격 부작용을 만들거나
  공유 상태를 변경할 수 없다.}
  \label{fig:architecture}
\end{figure}

\subsection{동결 입력과 근거 추적}

Worker 실행 전에 코디네이터는 논문 식별자, 논문 파일 신원, SHA-256과 크기,
근거 파일 신원과 hash 및 크기, 집계 input/evidence hash, agent version,
canonical task-prompt asset hash, schema hash를 논문별 manifest에 기록한다.
Native Worker instruction은 wrapper-package manifest로 별도 동결하고, 그
manifest hash는 batch metadata에 기록한다. 동결 corpus는 허용 경로를
제공하고, Root는 active lease를 manifest와 분리해 발급한다. 안전하지 않은
논문 식별자는 산출물이나 원장을 쓰기 전에 거부된다. Worker는 중앙 strength와
weakness에 대해 page, section, table, figure 또는 saved-result 위치와 함께
동결 신원 및 Root가 발급한 lease 신원을 반환해야 한다. 신원 불일치는 실행을
차단하며 targeted repair 경로로 들어갈 수 없다.

\subsection{정규 리뷰 구조}

초안에는 Summary, Strengths, Weaknesses, Questions, Ethics and Limitations,
Evidence Trace, score rationales가 포함된다. Soundness, Presentation,
Contribution, Significance, Originality, Overall Recommendation, Confidence,
constructive Comment는 서로 독립된 값이다. 결정론적 validator는 필수 필드,
정수 범위, 비어 있지 않은 prose, lease와 paper identity, evidence location,
금지된 filler text를 검사한다. 논문당 targeted repair는 schema와 calibration
오류를 합쳐 최대 한 번만 허용되며 동결 신원을 바꿀 수 없다.

\subsection{근거 우선 calibration}

JSON을 출력하기 전에 Worker는 비공개 claim map을 만들고, falsification
pass를 수행하며, 각 score를 독립적으로 anchor하고, claim, rationale, score,
confidence, comment 사이의 consistency를 검사한다. 결정론적 검증 뒤 Root는
복합 고위험 subset만 하나의 read-only verifier로 보낼 수 있다. 이 경로는
$\min(3,\lceil 0.3N\rceil)$편, 논문당 20초, T+15 이전, validated backlog가
최대 2일 때로 제한된다. 일반 초안은 fast path를 사용한다. 이는 실행 정책
제어이며 unseen paper에서의 리뷰 품질 향상으로 평가된 결과가 아니다.

\section{실행 방법}

운영자는 먼저 staged wrapper manifest를 검증하고, 두 Skill과 세 native-agent
정의를 project \texttt{.codex}에 설치한 뒤 새 Codex session을 시작한다. 호출
순서는 \texttt{\$auto-research}에서 Track 2-only frozen-paper 경로를 선택하고,
이어서 \texttt{\$ralphthon-track2-review-agent}를 실행하는 것이다. Root는 각
assignment를 동결하고, 논문별 drafting을 분배하며, deterministic validation과
제한된 risk gate를 실행하고, post 시도 전에 outbox와 clipboard 산출물을 쓴다.

포함된 \texttt{run\_batch.py}는 manifest 동결용 \texttt{BUILD}와 로컬 복구
검사용 mock-only \texttt{DRY-RUN}을 지원하며, output 생성 전에 \texttt{LIVE}를
거부한다. 승인된 live 시작 시점에는 Skill이 최대 45초의 read-only contract
discovery를 수행한다. Assignment, field, success marker 의미를 검증할 수 없으면
selector-free manual fallback으로 전환하되 동일한 validation 및 status-first
규칙을 유지한다.

\subsection{운영 체크리스트}

실제 실행 순서는 다음 네 단계다.
\begin{enumerate}
  \item \texttt{python3} \nolinkurl{scripts/install-track2-codex.py}
  \nolinkurl{--check-staging}으로 패키지를 확인한 뒤, 같은 명령을
  \nolinkurl{--install} 옵션으로 다시 실행하고 새 Codex session을 연다.
  \item \texttt{\$auto-research}를 호출해 Track 2-only 경로를 선택한다. 이어서
  다음 wrapper Skill을 호출한다.
  \begin{center}
    \texttt{\$ralphthon-track2-review-agent}
  \end{center}
  이때 \texttt{BUILD}, \texttt{DRY-RUN}, 또는 승인된 \texttt{LIVE} workflow
  중 하나를 명시한다.
  \item 로컬 준비에서는 \nolinkurl{run_batch.py}에 \nolinkurl{--mode}
  \texttt{BUILD}, \nolinkurl{--papers} \texttt{CORPUS},
  \nolinkurl{--output-dir} \texttt{DIR}을 넘긴다. 복구 리허설에서는 mode를
  \texttt{DRY-RUN}으로 바꾸고 \nolinkurl{--workers} \texttt{3}을 사용한다.
  생성된 manifest, ledger, outbox, clipboard, audit 결과를 보존한다.
  \item \texttt{LIVE}에서는 bundled CLI를 사용하지 않는다. Root가 제한된
  read-only discovery와 직렬화된 플랫폼 부작용을 담당하고, 관찰된 계약이
  불완전하면 \texttt{MANUAL\_PLATFORM.md}로 전환한다. 원격 완료를 확인했거나
  unresolved 상태를 명시적으로 기록했을 때만 실행을 종료한다.
\end{enumerate}

\section{신뢰성 모델}

코디네이터는 기본 high-water mark가 3인 제한된 queue, 논문당 하나의 active
lease, 원자적 snapshot 교체를 사용하는 append-only ledger, post intent용
idempotency key를 사용한다. timeout된 Worker는 제한된 retry 정책에 따라 한
번 다시 queue에 들어갈 수 있고, Root는 결정론적 lease를 유지한다. malformed
output은 공유 targeted repair 경로에 들어가기 전에 validator-reject 상태와
검증 오류를 기록한다.

Claim과 post timeout은 결과가 모호한 부작용이다. 따라서 코디네이터는 같은
동작을 다시 실행하기 전에 status를 조회한다. post 시도는 원격 상태가 verified
success 또는 안전한 retry로 해소될 때까지 \texttt{post\_unknown}으로 남는다.
유일한 live 완료 조건은 \texttt{posted\_verified == assigned\_count}이다.
Mock 완료는 mock 근거로만 보고하며 live-platform 성능으로 제시하지 않는다.

\section{평가 프로토콜}

평가는 quality 계약과 runtime 계약으로 나뉜다. Quality case는 candidate
generation 전에 issue type과 evidence location으로 seeded된다. 첫 번째 명시적
result sentence를 선택하는 고정 one-finding rule이 naive baseline이다.
결정론적 matching은 paper ID, finding type, normalized locator를 사용해 TP, FP,
FN을 계산한다. Finding text는 audit을 위해 보존하지만 match에는 포함되지
않으므로 이 검사는 semantic correctness를 평가하지 않는다. Acceptance
threshold는 결과 관찰 뒤에 고르는 대신 fixture와 함께 미리 저장된다.

Runtime 평가는 10개의 mock assignment와 세 번의 독립 반복을 사용한다.
assigned/verified-complete count, schema validity, duplicate post count, raw event
기반 latency를 기록한다. 별도 장애 시나리오는 malformed JSON, Worker timeout,
claim timeout, post timeout을 주입하고 원장을 한 번 in-process로 다시 연다.
각 시나리오는 verified recovery 또는 명시적 unresolved state를 보여야 하며,
조용한 유실은 통과가 아니다.

\section{합성 모의실험 결과}

표~\ref{tab:runtime}은 결정론적 synthetic fixture와 mock adapter만 보고한다.
세 번의 정상 실행에서 30개 assignment 모두 mock
\texttt{posted\_verified}에 도달했고, 30개 초안 모두 schema를 통과했으며,
duplicate post attempt와 duplicate idempotency key는 0이었다. 별도의 10편
실행은 malformed JSON output, Worker timeout, claim timeout, post timeout을
각각 한 번 주입하고 원장을 한 번 in-process로 다시 열었다. 주입된 조건은
각각 한 번 복구되었고, 10개 초안 모두 유효하고 mock-verified 상태에
도달했으며 duplicate post attempt는 0이었다. 완료 뒤 no-op rerun에서도
ledger와 outbox는 byte-identical 상태를 유지했다.

합성 모의실험 runtime의 OS-process resume 동작은 별도로 검사했다. 첫
프로세스는 mock-verified 4편 뒤 exit code 75로 종료했다. 새 프로세스는 같은
output directory와 input을 열고 10/10 mock verification, schema validity 10/10,
duplicate post/key 0으로 exit 0에 도달했다. 4편까지의 ledger prefix와 완료된
manifest/outbox byte도 보존됐다.

\begin{table}[H]
  \caption{결정론적 합성 모의실험 runtime 근거. Production claim 또는 post는
  발생하지 않았다.}
  \label{tab:runtime}
  \centering
  \small
  \begin{tabular}{@{}lrrr@{}}
    \toprule
    실행 & 검증 완료 & Schema & 중복 \\
    \midrule
    정상 실행 3회 & 30/30 & 30/30 & 0 \\
    장애 주입 & 10/10 & 10/10 & 0 \\
    두 프로세스 resume & 10/10 & 10/10 & 0 \\
    \bottomrule
  \end{tabular}
\end{table}

사전 동결된 seeded synthetic issue set에서 deterministic fixture-matching
runtime은 동결된 type/location target 20개 모두에 finding을 출력해 TP 20, FP 0,
FN 0, F1 1.0을 기록했다. 고정 one-finding baseline은 TP 10, FP 0, FN 10, F1
0.6667이었다. 동결된 rule에서 두 방법 모두 location accuracy 1.0이었다.
Finding prose를 채점하지 않으므로 이 결과는 scientific judgment나 semantic
review quality가 아니라 fixture type/location coverage를 검사한다.

\begin{table}[H]
  \caption{Seeded synthetic issue의 type/location coverage 평가.}
  \label{tab:quality}
  \centering
  \small
  \begin{tabular}{@{}lrrrr@{}}
    \toprule
    시스템 & TP & FP & FN & F1 \\
    \midrule
    Fixture-matching runtime & 20 & 0 & 0 & 1.0000 \\
    One-finding rule & 10 & 0 & 10 & 0.6667 \\
    \bottomrule
  \end{tabular}
\end{table}

Runtime과 quality 값은 다음 동결 기준 파일에서 가져왔다.
\begin{center}
\small\nolinkurl{evidence/mock-validation-final-20260712T1335KST/aggregate.json}
\end{center}
실제 두 프로세스 결과는 다음 파일에서 가져왔다.
\begin{center}
\small\nolinkurl{evidence/process-restart-proof/aggregate.json}
\end{center}
위 값은 저장 값의 직접 집계 또는 소수점 넷째 자리 반올림이다.

Post-build policy evaluator는 49개 regression test를 모두 통과했고,
installed/staged package의 정확한 일치와 정상 합성 모의실험의 30/30 유효,
duplicate-free 완료를 다시 확인했다. Native Worker instruction surface는
4,096-byte budget 이내인 2,224 bytes다. 이는 packaging과 fast-path 보존을
검증하지만 calibration의 과학적 효과를 입증하지 않는다.

\section{산출물과 감사 인계}

통과한 로컬 실행은 하나의 독립적인 evidence directory를 남긴다.
\texttt{manifests/}는 paper, evidence, prompt, schema, agent 신원을 묶고,
\texttt{outbox/}와 \texttt{clipboard/}는 post 시도 전에 검증된 review payload를
보존한다. \texttt{ledger.jsonl}, 원자적 state snapshot, \texttt{summary.json}은
상태 전이와 완료 수를 기록한다. 이 파일은 복구 입력이며 원격 제출 성공의
증거가 아니다.

재시작은 동일한 output directory와 동결 입력을 다시 사용한다. 이미 검증된
논문은 건너뛰지만, manifest 신원이 달라지면 덮어쓰지 않고 hard failure로
종료한다. Receipt audit은 assigned paper ID와 ledger의
\texttt{posted\_verified} 상태를 대조한다. Live manual lane에서 Root는 실제
플랫폼 success marker를 관찰한 뒤에만 이 상태를 기록한다. outbox 파일이나
로컬 mock 상태를 live receipt로 승격하지 않는다.

최종 live audit은 T+23:30에 시작하고 T+24:30에는 새로운 플랫폼 부작용을
중단한다. 인계 시 assigned, drafted, schema-valid, \texttt{post\_unknown},
\texttt{posted\_verified}, unresolved 수와 함께 \texttt{mock} 또는
\texttt{live} label을 명시한다. 따라서 부분 완료 batch도 성공으로 숨기지 않고
감사 가능한 unresolved 결과로 남는다.

\section{한계}

Mock run 통과만으로 live-platform throughput, submission success, unseen
research에서의 review correctness, human judgment와의 agreement를 입증할 수
없다. Runtime elapsed time은 in-process synthetic fixture를 반영하므로 live
speed claim에 사용하지 않는다. Worker-timeout 근거는 mock adapter
\texttt{TimeoutError}와 제한된 recovery를 사용하며, 실행 중 thread가 강제로
종료됐음을 입증하지 않는다.

구조적 한계도 있다. Evidence location은 claim이 도출된 위치를 보여줄 수
있지만 claim의 정확성을 보장하지 않는다. Originality 평가는 제출 논문과 제공
reference로 제한된다. 결정론적 schema validator는 malformed review를 거부할
수 있지만 scientific quality를 인증할 수 없다. Live adapter는 read-only
discovery로 실제 platform contract를 확인하기 전까지 정의되지 않으며, 본
보고서는 endpoint나 selector를 가정하지 않는다. Evidence-first와 verifier
정책은 prompt 및 orchestration 계약이며, unseen-paper 품질과 inter-reviewer
agreement에 미치는 효과는 측정하지 않았다.

\section{결론}

이 설계는 동결 근거 신원, 좁은 부작용 소유자, 제한된 논문별 Worker, 비공개
evidence-first calibration, 결정론적 validation, compound risk gate, status-first
reconciliation을 결합한다. 평가 계약은 fixture coverage를 semantic review
quality와 구분하고, mock evidence를 live outcome과 분리한다. 결정론적 mock
run은 동결 runtime, type/location coverage, 두 프로세스 resume 검사를
충족했지만 live-platform 동작과 unseen-paper review quality는 측정하지 않았다.

\begin{thebibliography}{2}

\bibitem{officialskill}
Ralphthon ICML organizers.
\newblock Auto Research Skill, Track 2-only frozen-paper path.
\newblock Repository commit
\nolinkurl{a9f4f2583648ef4ca54f980f951ae393d153473f}, 2026.

\bibitem{eventguide}
Ralphthon ICML organizers.
\newblock Track 2 rules, 2026.

\end{thebibliography}

\end{document}
